\chapter*{Introducción}
\addcontentsline{toc}{chapter}{Introducción}
El Renacimiento europeo de los siglos XV y XVI supuso una relectura activa de la Antigüedad clásica que transformó la cultura urbana, la vida política y las prácticas artísticas. Los humanistas impulsaron programas educativos basados en la retórica y la historia, al tiempo que los talleres artísticos y los gabinetes científicos incorporaron nuevos métodos de observación y registro.\cite{Kemp2008,AmesLewis2000} El auge de la imprenta multiplicó los circuitos de circulación de ideas, creando públicos lectores y una economía del libro que conectó ciudades distantes.\cite{Pettegree2010}

Esta monografía sigue la secuencia propuesta en el índice general para examinar el Renacimiento desde cinco ángulos complementarios: la relectura humanista de la Antigüedad, las artes visuales y la construcción de identidades, la relación entre ciencia y control social, la imprenta y la circulación del conocimiento, y los debates sobre género y sexualidad. Cada capítulo se apoya en bibliografía procedente de los volúmenes ubicados en la carpeta \texttt{libros\_descargados}, con referencias en estilo APA 7. El objetivo es ofrecer un panorama coherente que vincule contexto histórico, autores representativos y material gráfico extraído de las obras consultadas.
